\section{Reproducibility and claim traceability}
\label{app:reproducibility}

The supplementary source contains the frozen configuration, archived 32-row evidence table, aggregate summary, rerun summary, configuration digest, figure generator, and a shared-field comparison script. For anonymous review, the public location is represented as the \Repo; the camera-ready source should replace this placeholder with the permanent repository URL and commit.

\paragraph{Minimal reproduction.}
From the public source snapshot, install NumPy and pytest, set the repository root on \texttt{PYTHONPATH}, and run the three targeted test modules for the C1.1/C1.2 generator, evaluator and locked protocol. Then execute the locked runner with its explicit confirmation flag and a new empty output directory. Finally run \texttt{scripts/audit\_reproduction.py} from this package against the archived and fresh CSVs. The independent audit obtained 32/32 rows, 704/704 exact shared-field matches and maximum absolute error zero.

The archived CSV has 30 columns; the current public runner emits 22. The eight archive-only diagnostics are restart variance, composite incumbent feasibility, retained-state count, fallback use, $p=2$ near-optimal mass, effective support, and uniform/QAOA sample diversity. Because the schemas differ, the package does not claim byte identity. The richer archive remains the source for composite-feasibility and depth diagnostics; every field currently emitted by the runner is exactly reproduced.

\begingroup
\small
\begin{longtable}{@{}p{0.26\linewidth}p{0.34\linewidth}p{0.31\linewidth}@{}}
\caption{Claim-to-artifact traceability.}\label{tab:traceability}\\
\toprule
Claim / object & Public artifact or function & Verification route \\
\midrule
\endfirsthead
\toprule
Claim / object & Public artifact or function & Verification route \\
\midrule
\endhead
Generator and evaluator & \path{qops/c12_generator.py}; \texttt{generate\_instance}; \texttt{evaluate} & Exact formulas summarized in Appendix~\ref{app:benchmark}; targeted tests cover determinism and legality. \\
Classical controls & \path{qops/c12_classical.py}; coordinate greedy; cluster repair & Rerun all 32 frozen instances; compare gain, disagreement and feasibility fields. \\
Core and statevector & \path{qops/c12_quantum.py}; core extraction; retained-state QAOA & Recompute retained states, $p=1/p=2$ masses and matched draws. \\
Frozen matrix and inference & \path{qops/c12_locked.py}; locked config JSON & Explicit confirmation flag; configuration SHA-256; deterministic seeds. \\
Headline aggregates & Archived records CSV and summary JSON & 32 archived rows; fresh run matches 704/704 values shared with the current runner exactly. \\
Figures and tables & \path{scripts/generate_release_assets.py} and copied frozen CSV & Regenerate live-text SVG masters and table inputs; export PDF; run structural and numerical audits. \\
\bottomrule
\end{longtable}
\endgroup


\paragraph{Repository-wide health caveat.}
The targeted frozen experiment passes. A broad legacy test discovery also surfaced an unrelated sparse-tracking certificate failure, and a clean environment initially lacked pytest for a protocol module. These do not alter the C1.2 values, but they mean the snapshot should not be described as having an entirely green repository-wide test suite. The independent report records this distinction.
